\section{What prevents reliance on the present product for bank decisions}
\label{sec:legal-robustness-review}

The present evidence does not establish that LegalMath is suitable for autonomous
legal or compliance decisions in a bank. The code implements useful local
controls, but independent legal adjudication, performance on unseen regulatory
questions and the production identity, data and operational arrangements are
unfinished. The issue is not a residual percentage that can be inferred from
205 passing tests. Several necessary components of the intended use have not
yet been demonstrated.

This conclusion has consequences for product design. A proposed release needs
a named legal entity, regulated activity, source perimeter and decision that
the bank permits the software to influence. Review must distinguish the
authority governing that use from the technical evidence about the software.
Financial penalties, client harm and remediation duties cannot be estimated
from a software test count. Their assessment depends on the actual breach and
applicable legal and supervisory requirements.

\subsection{AI governance belongs in the regulatory perimeter}

The SFC's 12 November 2024 circular on generative AI language models addresses
licensed corporations' relevant use of such models, including third-party
components. It covers senior oversight, validation and ongoing management,
information security and provider risk. Its validation discussion includes
systems that combine a model with retrieval, prompts or other components.
The notification provisions concern the specified high-risk uses and the
existing notification requirements; they do not create a blanket approval
statement for every AI application \citep[paras.~6, 9--16, 20--31]{sfcgenai2024}.
An applicable-use assessment is therefore required before transferring these
expectations to a particular bank entity or workflow.

The accompanying appendix identifies failure modes relevant to this architecture,
including retrieval that misses necessary information and prompt injection.
It also discusses data and third-party risks. A citation-producing retrieval
system still needs review of its source support
\citep[appendix, paras.~9--12, 16--19]{sfcgenaiappendix2024}.

\begin{figure}[H]
\centering
\TeachingFlow{Entity, activity and AI use}{Applicable SFC or HKMA guidance}{Validation and operational responsibilities}
\caption{Governance obligations follow the actual use and institution.}
\label{fig:ai-governance-perimeter}
\end{figure}

The HKMA's 19 August 2024 circular specifically addresses consumer protection
in customer-facing generative-AI applications at authorized institutions. It
sets out governance, fairness, transparency and data-protection expectations,
including human control and review arrangements in the described circumstances
\citep[pp.~1--4]{hkmagenai2024}. LegalMath's internal authoring workbench and
a customer-facing explanation service are different uses. Enabling the latter
would require an additional applicability and control assessment. This cited
circular does not by itself determine the full risk-management requirements
for an internal bank system; its referenced guidance and other applicable
requirements need their own inventory.

For example, a synthetic control may correctly identify a failed wealth
condition but give a client an unsupported explanation about legal ineligibility.
That is a different harm from an incorrect internal Boolean calculation.
The product needs approved customer explanations, an accountable review route
and a record linking the explanation to the actual decision if such a use is
introduced. A technically accurate trace is a starting point for that design.

\subsection{Identity, access and evidence integrity need production controls}

The inspected local implementation registers synthetic identities and their
roles, and binds reviews to exact content. This is sufficient to exercise
separation of author, meaning reviewer and engineering reviewer in tests.
It does not establish an employee's current delegated authority, employment
status, multi-factor authentication or revocation. A bank identity integration
must revoke a departing reviewer's future authority while preserving the
historical record of decisions genuinely authorised at the time.

Access requires a separate design. Authentication of a registered local user
does not establish that the user may inspect every client's evidence or every
legal entity's case. Production permissions must apply to source packets,
candidate reports, exports and factual evidence, including indirect retrieval
through linked records. A test should attempt to retrieve a second entity's
case through every relevant interface and verify the intended denial. The
current local API is not evidence that this separation has been implemented.

\begin{figure}[H]
\centering
\TeachingCompare{Who is the caller?}{Which case may they inspect?}{Which decision may they approve?}
\caption{Authentication, access and decision authority require separate answers.}
\label{fig:identity-access-authority}
\end{figure}

The audit chain is stored in the same locally administered environment as the
database it records. Its hashes can expose changed records when the expected
chain remains trusted. An administrator who can replace both records and the
chain can defeat that comparison. The bank design therefore needs an independent
integrity reference, such as appropriately controlled signed checkpoints or
external immutable retention, with reviewed key custody and recovery. This is
a design requirement arising from the trust model, not a claim that one named
technology is mandated by every regulator.

Consider a synthetic incident in which a privileged operator replaces an approval
and recomputes the later hashes. A verifier that trusts only the rewritten
database can find an internally consistent history. An independently retained
earlier checkpoint would instead expose the divergence. Testing accidental blob
corruption does not exercise this privileged-rewrite scenario.

\subsection{Abstention needs an operational destination}

An unresolved interpretation is useful only if the affected process handles it.
Every material issue needs an owner, an evidence request, a permitted interim
action and a review deadline appropriate to the bank's workflow. The system
must detect an unowned issue or an overdue queue rather than allowing a blocked
screen to become routine ignored noise. An override needs defined authority,
scope and expiry; it must not silently amend the underlying regulatory meaning.

\begin{figure}[H]
\centering
\TeachingFlow{Unresolved classification}{Named review owner and interim action}{Resolution or explicit continued restriction}
\caption{Abstention transfers a specific decision; it does not finish the operating process.}
\label{fig:abstention-owner}
\end{figure}

The pure Java evaluator also needs a production release-withdrawal protocol.
An offline library can continue executing after its central approval is retired
unless the host checks the applicable release policy. The bank must define
how revocation reaches hosts, the permitted age of cached authority, behavior
during disconnection and treatment of in-flight transactions. A release revoked
for wrong meaning raises a different fallback question from a model-provider
outage while an approved rule remains applicable.

For a synthetic race, an order is assessed just before a release is suspended,
then reaches commit after suspension. The host must apply its reviewed policy
to that transition and record the version used. Merely retaining the JAR hash
does not decide whether the order may complete. The existing in-memory host
example motivates this test but does not validate a bank's distributed order
and deployment infrastructure.

\subsection{Confidentiality and incidents cross the model boundary}

Public regulatory documents and synthetic examples permit development without
client data. Production work can introduce personal information, confidential
advice and privileged legal analysis. The bank must identify which material can
leave its environment, for what purpose, under which provider terms, and with
what retention, deletion and access arrangements. Privilege and cross-border
handling require qualified assessment for the actual arrangement. Encrypting
transport alone does not answer those questions.

A source parser and a model provider also present different attack surfaces.
The parser needs operating-system isolation and resource bounds for hostile
files. The provider needs a narrow task interface and cannot inherit approval,
network or database privileges from text it reads. The complete application
must be tested with malicious source instructions, altered attachments and
cross-case references. The focused security, release-binding and host checks
rerun for this review passed 18 tests; that result covers those cases only.

\begin{figure}[H]
\centering
\TeachingFlow{Preserve failing source and decision}{Contain and trace affected uses}{Review remediation and notification duties}
\caption{Incident response requires evidence, operational action and legal assessment.}
\label{fig:incident-legal-response}
\end{figure}

Finally, an incident procedure must assign authority for suspension, client
remediation, regulatory-notification assessment and preservation of relevant
records. The product should expose affected decisions and their dependencies
without deciding the bank's external communications. Rehearsing a missing
exception, a compromised approval and a stale release tests three different
response paths. These exercises, independent legal review and a measured shadow
study are necessary next evidence for a bounded deployment decision. The current
document and local tests cannot substitute for them.
